\chapter{Analisi del grafo bipartito: vedere la struttura}
\label{ch:metodo_grafo}

\epigraph{``A network is more than the sum of its nodes.''}{Mark
E. J. Newman, \emph{Networks: An Introduction}~\parencite{newman2010}}

\lettrine[lines=3]{O}{gni} scuola pu\`o essere vista come un
\emph{grafo bipartito} che ha per nodi le classi e i docenti,
con un arco fra un docente e una classe se il docente
insegna in quella classe (eventualmente pesato dal numero di
ore). Questo grafo non \`e un'astrazione gratuita: ne
emergono \emph{misure quantitative} che dicono molto sulla
struttura organizzativa della scuola e sulla difficolt\`a del
problema di scheduling.

In $\pi$Tantum la pagina Diagnostica espone tre misure
classiche della network analysis:
\emph{modularit\`a}\index{modularit\`a},
\emph{betweenness centrality}\index{betweenness} e
\emph{densit\`a}\index{densit\`a (grafo)}. Ciascuna ha un
significato preciso e un suggerimento pratico associato.

\section{Modularit\`a: quanto \`e ``a cluster'' la scuola}

\begin{storiabox}
La modularit\`a \`e una misura introdotta da Newman e Girvan
nel 2004~\parencite{newman2004} per quantificare la \emph{community
structure} di un grafo: \`e alta quando il grafo si divide
naturalmente in gruppi (comunit\`a) ben separati, \`e bassa o
negativa quando il grafo \`e omogeneo. \`E diventata una delle
misure pi\`u usate nella scienza delle reti.
\end{storiabox}

\begin{defbox}
Per un grafo $G$ con $m$ archi e una partizione dei nodi in
comunit\`a $C_1, \ldots, C_k$, la modularit\`a \`e:
\[
  Q = \frac{1}{2m} \sum_{i,j} \left[ A_{ij} - \frac{k_i k_j}{2m}
                            \right] \delta(c_i, c_j)
\]
dove $A_{ij}$ \`e la matrice di adiacenza, $k_i$ \`e il grado
del nodo $i$, $\delta(c_i, c_j) = 1$ se $i$ e $j$ sono nella
stessa comunit\`a, 0 altrimenti.
\end{defbox}

\textbf{Come si legge}: per ogni coppia di nodi nella stessa
comunit\`a, si confronta il \emph{numero effettivo} di archi
fra loro ($A_{ij}$) con il \emph{numero atteso} se gli archi
fossero distribuiti casualmente preservando i gradi
($k_i k_j / 2m$). Se i nodi della stessa comunit\`a sono
\emph{pi\`u connessi} di quanto ci si aspetterebbe a caso, $Q$
\`e positiva. Valori tipici: $Q < 0.3$ (struttura debole),
$0.3 < Q < 0.5$ (struttura moderata), $Q > 0.5$ (struttura
marcata).

\begin{esempiobox}
La scuola di fig.~\ref{fig:spettrale-cluster}
(cap.~\ref{ch:metodo_spettrale}) ha modularit\`a $Q \approx
0.42$ (struttura moderata: due cluster ben separati ma con
un ponte). Una scuola con un solo indirizzo e docenti che
insegnano un po' ovunque ha tipicamente $Q < 0.2$.
\end{esempiobox}

\textbf{Suggerimento pratico}: se $Q > 0.4$, la
decomposizione spettrale (cap.~\ref{ch:metodo_spettrale})
funzioner\`a bene; la pipeline pu\`o essere accelerata
risolvendo cluster indipendenti. Se $Q < 0.2$, la
decomposizione produrr\`a cluster artificiali e conviene
risolvere monoliticamente.

\section{Betweenness centrality: i ``ponti'' del sistema}

\begin{storiabox}
La \emph{betweenness} fu introdotta da Linton Freeman nel
1977~\parencite{freeman1977} nella sociologia dei piccoli
gruppi: misura quanto un individuo \`e ``intermediario'' nelle
comunicazioni del gruppo. Si \`e poi diffusa in tantissimi
campi (citation networks, Internet, biology).
\end{storiabox}

\begin{defbox}
Per un nodo $v$, la \emph{betweenness centrality} \`e:
\[
  C_B(v) = \sum_{s \neq v \neq t} \frac{\sigma_{st}(v)}{\sigma_{st}}
\]
dove $\sigma_{st}$ \`e il numero di shortest path da $s$ a
$t$, $\sigma_{st}(v)$ il numero di quelli che passano per
$v$.
\end{defbox}

\textbf{Come si legge}: $C_B(v)$ \`e alta se $v$ \`e ``in
mezzo'' a molti percorsi del grafo. Un docente con
betweenness alta \`e quello che, se si toglie, ``sconnette''
parecchio il sistema. Tipicamente sono docenti che
insegnano in pi\`u indirizzi diversi, fungendo da ponte fra
cluster.

\textbf{Suggerimento pratico}: i top-$k$ docenti per
betweenness sono quelli su cui il timetabling \`e pi\`u
fragile. Una loro indisponibilit\`a aggiuntiva ha alta
probabilit\`a di rendere il problema infeasible. Vale la
pena tenerli sotto controllo durante l'anno.

\section{Densit\`a: quanto \`e ``pieno'' il grafo}

\begin{defbox}
La \emph{densit\`a} di un grafo bipartito con $n_C$ classi e
$n_D$ docenti, che ha $|E|$ archi, \`e:
\[
  \rho = \frac{|E|}{n_C \cdot n_D}
\]
\end{defbox}

Vale fra 0 (nessun arco) e 1 (tutti i docenti insegnano in
tutte le classi). Per scuole italiane reali, $\rho$ \`e
tipicamente $0.05$--$0.15$.

\textbf{Suggerimento pratico}: $\rho$ alta significa molti
docenti che insegnano in molte classi (es. ITIS dove i
docenti coprono pi\`u indirizzi). $\pi$Tantum tipicamente
risolve in pi\`u tempo le scuole con $\rho$ alta perch\'e i
vincoli fra docenti diversi sono pi\`u intrecciati. Valori
$\rho > 0.3$ sono considerati ``densi'' e suggeriscono di
\emph{non} usare la decomposizione spettrale (vedi sopra).

\section{Visualizzazione e interpretazione}

In Diagnostica $\pi$Tantum mostra:
\begin{itemize}
\item Una \emph{matrice riassuntiva} con $Q$, $\rho$,
  numero di componenti connesse, top-5 docenti per
  betweenness.
\item Un \emph{istogramma} dei gradi (quante classi per
  docente; quanti docenti per classe).
\item Un \emph{heatmap} della matrice di adiacenza con
  righe ordinate per cluster spettrale: a colpo d'occhio
  si vedono i cluster e i ponti.
\end{itemize}

\begin{casostudiobox}
\textbf{Liceo ``Manzoni'', 24 classi}.
$Q = 0.51$, $\rho = 0.09$, top docente per betweenness:
prof Esposito (insegna A026 in liceo classico e in liceo
scientifico, fa da ponte fra i due). Modularit\`a alta
$\to$ decomposizione spettrale produce 3 cluster (classico,
scientifico, linguistico) ben separati. Esposito \`e l'unico
ponte significativo; basta tenerlo come ``shared'' fra due
cluster con vincolo lagrangiano. Pipeline completa in
4~min.
\end{casostudiobox}

\section{Connessioni}

\begin{connessionebox}
\begin{itemize}
\item \textbf{Decomposizione spettrale}
  (cap.~\ref{ch:metodo_spettrale}): la modularit\`a $Q$
  predice la qualit\`a della decomposizione.
\item \textbf{Hall pre-check} (cap.~\ref{ch:metodo_hall}):
  $\rho$ alta in pratica significa molti vincoli di
  matching da verificare; Hall richiede pi\`u tempo.
\item \textbf{Monte Carlo} (cap.~\ref{ch:metodo_montecarlo}):
  i docenti con alta betweenness sono i candidati naturali
  per le perturbazioni del Monte Carlo (testare cosa
  succede se diventano indisponibili).
\end{itemize}
\end{connessionebox}

\begin{biblioparvabox}
\begin{itemize}
\item \cite{newman2004}: il paper originale sulla
  modularit\`a (Newman \& Girvan).
\item \cite{freeman1977}: il paper originale sulla
  betweenness.
\item \cite{newman2010}: \emph{Networks: An
  Introduction}, manuale moderno e accessibile per la
  network analysis.
\end{itemize}
\end{biblioparvabox}
